\begin{onehalfspace}

Questo lavoro è partito da una domanda operativa: quanto dell'output
rumoroso di Shor può essere recuperato a valle, e quando diventa invece
necessario correggere l'errore durante il calcolo? La risposta è una
delimitazione, non una promessa generale. Su $N=15$ una semplice ricerca
multi-candidato usa meglio l'informazione già presente; un classificatore
sullo stesso output non aggiunge valore. Per andare oltre occorre
spostarsi sulla sindrome e sulla correzione d'errore.

% -----------------------------------------------
\section{Le Domande di Ricerca}
\label{sec:risposte_dr}

\paragraph{DR1 --- Come cambia Shor al variare dei canali di rumore?}
Nel modello uniforme N=15 la struttura QPE si attenua all'aumentare di
$\lambda_{2q}$: la frazione efficace sui quattro picchi passa da 0.804 a
$\lambda_{2q}=10^{-3}$ a 0.042 a $10^{-1}$. Il proxy di nessun evento
Pauli sulle 224 CX decade molto più rapidamente, da 0.811 a
$2.65\times10^{-10}$, e non predice il successo dell'algoritmo. Gli
sweep di $\lambda_{1q}$, $T_1/T_2$, readout e shot non mostrano un
andamento monotono di M1 con 30 repliche; non è quindi lecito ordinarne
causalmente l'importanza da questi soli dati. TOP-4 resta a una
iterazione in tutti i punti salvo $\lambda_{2q}=0.1$, dove sale a 1.07.

\paragraph{DR2 --- Come si estrae una sindrome senza misurare lo stato
logico?}
Il codice a ripetizione misura parità, non i singoli qubit. Gli
stabilizzatori $Z_1Z_2$ e $Z_2Z_3$ rivelano quale bit è stato alterato
senza rivelare l'ampiezza logica. La curva Monte Carlo coincide con
$p_L=3p^2-2p^3$ entro l'incertezza; il limite resta la protezione di un
solo tipo di errore per volta.

\paragraph{DR3 --- Quali vantaggi e limiti offre Steane $[[7,1,3]]$?}
La struttura CSS corregge un errore Pauli arbitrario su un singolo qubit
e la corrispondenza sindrome--errore è stata verificata sui 21 casi. La
pendenza log--log misurata, 1.94, è compatibile con la soppressione
quadratica attesa da un codice di distanza 3. La distanza fissa e
l'estrazione di sindrome non fault-tolerant impediscono però di
interpretarlo come soluzione scalabile.

\paragraph{DR4 --- Il surface code mostra un comportamento di soglia?}
Sì, nel memory experiment e nel modello circuit-level adottati. Le curve
per $d=3,5,7,9$ si incrociano attorno allo $0.9\%$ in base Z e allo
$0.7\%$ in base X; sotto la regione di incrocio aumentare la distanza
sopprime il logical error rate, sopra la amplifica. Il parametro fisico
$p$ di questa simulazione non coincide con $\lambda_{2q}$ della campagna
Shor, quindi il confronto fra i due resta qualitativo.

\paragraph{DR5 --- Quale errore logico è compatibile con Shor?}
Il modello fenomenologico mostra come la probabilità per singola misura
decada al crescere della probabilità di Pauli non-identità per gate
logico. A $p_L=0$ il limite dell'istanza è circa 0.75, perché una parte
degli esiti QPE è intrinsecamente non utile; ripetere l'algoritmo aumenta
la probabilità cumulativa. La curva è un'analisi di sensibilità, non una
previsione end-to-end: il logical error rate di un memory experiment non
può essere assegnato direttamente a ogni gate di Shor senza specificare
le operazioni fault-tolerant e i rispettivi canali residui.

\paragraph{DR6 --- Quando un decoder appreso può superare quello
analitico?}
Nei dati della tesi il vantaggio compare quando il rumore contiene
correlazioni assenti dal grafo del matching. Il decoder ibrido riduce
l'errore rispetto a \ac{MWPM} a distanza 3, ma il beneficio si restringe
con la distanza e scompare a distanza 7 per la rete densa adottata. Il
risultato non sostiene la sostituzione generalizzata dei decoder
analitici: indica una nicchia in cui un modello residuale può sfruttare
informazione strutturale non rappresentata dal metodo base.

% -----------------------------------------------
\section{Verifica delle Ipotesi}
\label{sec:verifica_ipotesi}

\begin{table}[H]
\centering
\small
\caption[Esito delle ipotesi]{Esito dopo la rigenerazione schema~2.}
\label{tab:esito_ipotesi}
\begin{tabular}{p{4.2cm}p{2.5cm}p{6.6cm}}
\toprule
\textbf{Ipotesi} & \textbf{Esito} & \textbf{Evidenza} \\
\midrule
Il classificatore riduce le iterazioni
& smentita
& M2 non migliora M1; TOP-4 è migliore di M2 in UC1 e UC2
($p_{\mathrm{Holm}}=0.0033/0.0096$). \\
\addlinespace
F1 $>0.85$ e AUC $>0.90$
& confermata, ma insufficiente
& SVM: F1 $0.919/0.923$, AUC $0.965/0.958$ sul test; la metrica
misura TOP-1, non l'utilità operativa. \\
\addlinespace
Il vantaggio si mantiene crescendo $N$
& non verificata
& $N=21/35$ sono validati idealmente, ma esclusi dal rumore per costo;
l'ordine $r$ deve essere controllato insieme a $N$. \\
\addlinespace
Soppressione quadratica nei codici di distanza 3
& confermata nel modello
& ripetizione coerente con $3p^2-2p^3$; Steane con pendenza 1.94. \\
\addlinespace
Comportamento di soglia del surface code
& confermato nel modello
& incrocio delle curve $d=3,5,7,9$ e inversione sopra soglia. \\
\addlinespace
Trasferimento diretto di $p_L$ a Shor
& non dimostrato
& la curva logica è fenomenologica; manca la mappatura per operazione
fault-tolerant. \\
\bottomrule
\end{tabular}
\end{table}

L'audit del classificatore è metodologicamente centrale. TOP-1 produce
classi addestrabili, ma chiede soltanto se il candidato dominante porta
ai fattori. TOP-16 produce invece $2\,000/2\,000$ positivi in entrambi
gli scenari. Il conto combinatorio spiega il fenomeno: 63 dei 256 esiti
sono utili e la probabilità di non includerne alcuno scegliendo sedici
celle uniformi è appena
$\binom{193}{16}/\binom{256}{16}=0.9275\%$. Una metrica elevata non
garantisce dunque che l'etichetta rappresenti un problema utile.

% -----------------------------------------------
\section{Contributi del Lavoro}

\textbf{Una demo scientificamente separata in ideale e rumoroso.}
Il percorso ideale espone circuito, stato e QPE senza rumore; il percorso
rumoroso rende attivabili separatamente depolarizzazione 1Q/2Q,
$T_1/T_2$, readout e perturbazione coerente. N=21/35 sono validati
idealmente, mentre l'esperimento pubblico rumoroso è limitato a N=15.

\textbf{Un'ablazione causalmente interpretabile.}
M1, TOP-4 e M2 ricevono lo stesso istogramma. TOP-4 riduce $\bar M$ da
1.37 a 1.00 nello scenario moderato e da 1.50 a 1.00 nello stress; il
classificatore peggiora la propria ablazione. Seed, tie-break,
sentinelle, revisioni e manifest sono salvati negli artefatti.

\textbf{Una validazione funzionale dell'aritmetica.}
Il moltiplicatore $\times7\bmod15$ è coperto da truth table e il percorso
Beauregard da proprietà esaustive e confronti QPE per N=21/35. Questo
controllo ha mostrato perché i soli picchi o il solo ordine non bastano a
validare l'operatore implementato.

\textbf{Una progressione QEC con limiti dichiarati.}
Ripetizione, Steane, surface code e decodifica appresa sono collegati da
metriche esplicite, ma i rispettivi modelli di rumore non vengono
presentati come direttamente commensurabili. Lo Shor logico chiude la
progressione come analisi fenomenologica di sensibilità.

% -----------------------------------------------
\section{Limiti}
\label{sec:limiti_finali}

La campagna Shor usa un simulatore e un rumore uniforme, indipendente e
stazionario; non include topologia, routing, crosstalk, leakage o deriva
di calibrazione. N=15 è un'istanza didattica e TOP-K beneficia di una
post-elaborazione particolarmente permissiva. Gli sweep hanno 30
repliche e i test non sono corretti globalmente fra le otto famiglie.
N=21/35 non sono simulati sotto rumore.

Nella parte QEC, Steane non usa un'estrazione fault-tolerant; i memory
experiment e il modello per-gate di Shor non condividono la stessa
definizione di $p_L$; il decoder appreso è una rete densa e il vantaggio
non persiste alle distanze maggiori testate. Le stime di risorse
crittografiche devono quindi provenire da architetture fault-tolerant
complete, non da un'estrapolazione delle curve N=15.

% -----------------------------------------------
\section{Il Criterio che Unifica il Lavoro}

In entrambe le fasi, il complemento semplice è utile quando sfrutta
informazione che il metodo base lascia inutilizzata senza modificare il
processo che la genera. TOP-4 prova altri candidati già presenti
nell'istogramma; il decoder residuale interviene soltanto sui casi in cui
il matching ha una debolezza rappresentazionale. Un modello più
complesso non produce valore se replica una regola semplice, se le
etichette sono degeneri o se il costo introdotto compensa il beneficio.

Applicato a Shor, questo criterio separa nettamente post-processing e
correzione: TOP-K può recuperare un picco utile già misurato, ma non può
ricostruire informazione fisicamente distrutta. Rendere Shor scalabile
richiede operazioni logiche fault-tolerant; l'apprendimento può contribuire
alla decodifica solo dove dispone di una struttura che il decoder
analitico non rappresenta.

\end{onehalfspace}
